\documentclass[12pt]{article}
\usepackage[T1]{fontenc}
\usepackage[T1]{polski}
\usepackage[utf8]{inputenc}
\newcommand{\BibTeX}{{\sc Bib}\TeX} 
\usepackage{graphicx}
\usepackage{amsfonts}
\usepackage{amsmath}
\usepackage{latexsym}
\usepackage{float}

\setlength{\textheight}{21cm}

\title{{\bf Zadanie nr 1 - Generacja sygnału i szumu}\linebreak
Cyfrowe Przetwarzanie Sygnałów}
\author{Dawid Jakubik, 224307 \and Hubert Gawłowski, 224298}
\date{22.04.2021}

\begin{document}
\clearpage\maketitle
\thispagestyle{empty}
\newpage
\setcounter{page}{1}
\section{Cel zadania}

Celem zadania było zapoznanie się z procesem konwersji A/C oraz C/A sygnałów jednocześnie poznając kilka metod na przeprowadzenie obu procesów jak i zaiplementowanie ich w aplikacji. 

\section{Wstęp teoretyczny}

Sygnały, które podlegaja kwantyzacji oraz rekonstrukcji generowane są na podstawie wzorów znajdujących się w instrukcji do zadania nr 1 \cite{instrukcja1}. W instrukcji \cite{instrukcji2} znajdują się wzory oraz opisy metod konwersji A/C oraz C/A użyte w celu kwantyzacji, rekonstrukcji oraz do obliczenia błędu średniokwadratowego, stosunku sygnał-szum, szczytowego stosunku sygnał-szum oraz maksymalnej różnicy.


\section{Eksperymenty i wyniki}

%%%%%%%%%%%%%%%%%%%%%%%%%%%%%%%%%%%%%%%%%%%%%%%%%%%%%%%%%%%%%%%%%%%%%%%%%%%%%%%%%%%%%%%%%%%%%%%%%%%%%%%%%%%%%%%%%
% PODROZDZIA£ PT. EKSPERYMENT NR 1 
%%%%%%%%%%%%%%%%%%%%%%%%%%%%%%%%%%%%%%%%%%%%%%%%%%%%%%%%%%%%%%%%%%%%%%%%%%%%%%%%%%%%%%%%%%%%%%%%%%%%%%%%%%%%%%%%%
W każdym z eksperymantów jako sygnał ciągły posłużył nam sygnał sinusoidalny o następujących parametrach: 
\begin{itemize}
	\item Amplituda: 4
	\item Czas początkowy: 0s
	\item Czas trwania sygnału: 2s
	\item Okres: 1s
\end{itemize}

\subsection{Eksperyment nr 1 Kwantyzacja równomierna z obcięciem}

Celem eksperymentu było przeprowadzenie kwantyzacji sygnału ciągłego z obcięciem.

\subsubsection{Założenia}

\begin{itemize}
    \item Częstotliwość próbkowania:
    \item Liczba bitów zapisu poziomów kwantyzacji:
\end{itemize}
\subsubsection{Przebieg}
Wygenerowany sygnał oraz jego skwantyzowana postać prezentują sie następująco:
\begin{figure}[H]
    \centering
    %\includegraphics{cps_kwantyzacja_z_obcieciem.jpg}
	\includegraphics[width=\linewidth]{tmp.jpg}
    \caption{Kwantyzacja równomierna z obcięciem sygnału sinusoidalnego}
    \label{wykres dla eksperymentu 1}
\end{figure}

\subsubsection{Rezultat}
Obliczone miary dla procesu kwantyzacji:
\begin{figure}[H]
    \centering
    %\includegraphics{cps_kwantyzacja_z_obcieciem_miary.jpg}
	\includegraphics[width=\linewidth]{tmp.jpg}
    \caption{Miary obliczone dla kwantyzacji równomiernej z obcięciem sygnału sinusoidalnego}
    \label{Wartości dla eksperymentu 1}
\end{figure}


%%%%%%%%%%%%%%%%%%%%%%%%%%%%%%%%%%%%%%%%%%%%%%%%%%%%%%%%%%%%%%%%%%%%%%%%%%%%%%%%%%%%%%%%%%%%%%%%%%%%%%%%%%%%%%%%%

%%%%%%%%%%%%%%%%%%%%%%%%%%%%%%%%%%%%%%%%%%%%%%%%%%%%%%%%%%%%%%%%%%%%%%%%%%%%%%%%%%%%%%%%%%%%%%%%%%%%%%%%%%%%%%%%%

\newpage
\subsection{Eksperyment nr 2 : Kwantyzacja równomierna z zaokrąglaniem}

Celem eksperymentu było przeprowadzenie kwantyzacji sygnału ciągłego z zaokrągleniem.

\subsubsection{Założenia}

\begin{itemize}
	\item Częstotliwość próbkowania:
	\item Liczba bitów zapisu poziomów kwantyzacji:
\end{itemize}
\subsubsection{Przebieg}
Wygenerowany sygnał oraz jego skwantyzowana postać prezentują sie następująco:
\begin{figure}[H]
	\centering
	%\includegraphics{cps_kwantyzacja_z_zaokrągleniem.jpg}
	\includegraphics[width=\linewidth]{tmp.jpg}
	\caption{Kwantyzacja równomierna z zaokrągleniem sygnału sinusoidalnego}
	\label{wykres dla eksperymentu 2}
\end{figure}

\subsubsection{Rezultat}
Obliczone miary dla procesu kwantyzacji:
\begin{figure}[H]
	\centering
	%\includegraphics{cps_kwantyzacja_z_zaokrągleniem_miary.jpg}
	\includegraphics[width=\linewidth]{tmp.jpg}
	\caption{Miary obliczone dla kwantyzacji równomiernej z zaokrągleniem sygnału sinusoidalnego}
	\label{Wartości dla eksperymentu 2}
\end{figure}


%%%%%%%%%%%%%%%%%%%%%%%%%%%%%%%%%%%%%%%%%%%%%%%%%%%%%%%%%%%%%%%%%%%%%%%%%%%%%%%%%%%%%%%%%%%%%%%%%%%%%%%%%%%%%%%%%

%%%%%%%%%%%%%%%%%%%%%%%%%%%%%%%%%%%%%%%%%%%%%%%%%%%%%%%%%%%%%%%%%%%%%%%%%%%%%%%%%%%%%%%%%%%%%%%%%%%%%%%%%%%%%%%%%
\newpage
\subsection{Eksperyment nr 3: Ekstrapolacja zerowego rzędu}


\subsubsection{Założenia}
Eksperyment nr 3 polegał na ekstrapolacji zerowego rzędu spróbkowanego sygnału ciągłego.
\subsubsection{Przebieg}
Wygenerowany sygnał prezentuje się następująco:
\begin{figure}[H]
    \centering
    %\includegraphics{cps_ekstrapolacja_0.jpg}
	\includegraphics[width=\linewidth]{tmp.jpg}
    \caption{Wykres Ekstrapolacji zerowego rzędu sygnału sinusoidalnego}
    \label{wykres dla eksperymentu 3}
\end{figure}



\subsubsection{Rezultat}
Obliczone parametry dla otrzymanego w rezuntacie sygnału ciągłego:
\begin{figure}[H]
    \centering
    %\includegraphics{cps_ekstrapolacja_0_miary.jpg}
	\includegraphics[width=\linewidth]{tmp.jpg}
    \caption{Oblicznone miary dla Ekstrapolacji zerowego rzędu sygnału sinusoidalnego}
    \label{wartości dla eksperymentu 3}
\end{figure}


%%%%%%%%%%%%%%%%%%%%%%%%%%%%%%%%%%%%%%%%%%%%%%%%%%%%%%%%%%%%%%%%%%%%%%%%%%%%%%%%%%%%%%%%%%%%%%%%%%%%%%%%%%%%%%%%%

%%%%%%%%%%%%%%%%%%%%%%%%%%%%%%%%%%%%%%%%%%%%%%%%%%%%%%%%%%%%%%%%%%%%%%%%%%%%%%%%%%%%%%%%%%%%%%%%%%%%%%%%%%%%%%%%%
\newpage
\subsection{Eksperyment nr 4: Interpolacja pierwszego rzędu}

\subsubsection{Założenia}
Eksperyment nr 4 polegał naiInterpolacji pierwszego rzędu spróbkowanego sygnału ciągłego.
\subsubsection{Przebieg}
Wygenerowany sygnał prezentuje się następująco:
\begin{figure}[H]
	\centering
	%\includegraphics{cps_interpolacja_1.jpg}
	\includegraphics[width=\linewidth]{tmp.jpg}
	\caption{Wykres Interpolacji pierwszego rzędu sygnału sinusoidalnego}
	\label{wykres dla eksperymentu 4}
\end{figure}



\subsubsection{Rezultat}
Obliczone parametry dla otrzymanego w rezuntacie sygnału ciągłego:
\begin{figure}[H]
	\centering
	%\includegraphics{cps_interpolacja_1_miary.jpg}
	\includegraphics[width=\linewidth]{tmp.jpg}
	\caption{Oblicznone miary dla Interpolacja pierwszego rzędu sygnału sinusoidalnego}
	\label{wartości dla eksperymentu 4}
\end{figure}

%%%%%%%%%%%%%%%%%%%%%%%%%%%%%%%%%%%%%%%%%%%%%%%%%%%%%%%%%%%%%%%%%%%%%%%%%%%%%%%%%%%%%%%%%%%%%%%%%%%%%%%%%%%%%%%%%

%%%%%%%%%%%%%%%%%%%%%%%%%%%%%%%%%%%%%%%%%%%%%%%%%%%%%%%%%%%%%%%%%%%%%%%%%%%%%%%%%%%%%%%%%%%%%%%%%%%%%%%%%%%%%%%%%
\newpage
\subsection{Eksperyment nr 5: Rekonstrukcja w oparciu o funkcję sinc }
\subsubsection{Założenia}
Eksperyment nr 3 polegał na rekonstrukcji w oparciu o funkcję sinc spróbkowanego sygnału ciągłego.
\subsubsection{Przebieg}
Wygenerowany sygnał prezentuje się następująco:
\begin{figure}[H]
	\centering
	%\includegraphics{cps_rekonstrukcja_sinc.jpg}
	\includegraphics[width=\linewidth]{tmp.jpg}
	\caption{Wykres rekonstrukcji w oparciu o funkcję sinc sygnału sinusoidalnego}
	\label{wykres dla eksperymentu 5}
\end{figure}



\subsubsection{Rezultat}
Obliczone parametry dla otrzymanego w rezuntacie sygnału ciągłego:
\begin{figure}[H]
	\centering
	%\includegraphics{cps_rekonstrukcja_sincmiary.jpg}
	\includegraphics[width=\linewidth]{tmp.jpg}
	\caption{Oblicznone miary dla rekonstrukcji w oparciu o funkcję sinc sygnału sinusoidalnego}
	\label{wartości dla eksperymentu 5}
\end{figure}
%%%%%%%%%%%%%%%%%%%%%%%%%%%%%%%%%%%%%%%%%%%%%%%%%%%%%%%%%%%%%%%%%%%%%%%%%%%%%%%%%%%%%%%%%%%%%%%%%%%%%%%%%%%%%%%%%

%%%%%%%%%%%%%%%%%%%%%%%%%%%%%%%%%%%%%%%%%%%%%%%%%%%%%%%%%%%%%%%%%%%%%%%%%%%%%%%%%%%%%%%%%%%%%%%%%%%%%%%%%%%%%%%%%
\newpage
\subsection{Eksperyment nr 6: Zjawisko aliasingu }
\subsubsection{Założenia}
Eksperyment nr 6 polegał na rekonstrukcji w oparciu o funkcję sinc spróbkowanego sygnału ciągłego z tak dobranymi parametrami by uzyskac zjawisko aliasingu.
\subsubsection{Przebieg}
Wygenerowane sygnały sinusoidalne o częstotliwościach 100Hz, 220Hz oraz 440Hz  zostały spróbkowane kolejno z częstotliwościami  1 000Hz, 44 100Hz oraz22 050Hz a następnie sygnały zostały zrekonstruowane za pomocą finkcji sinc.

\subsubsection{Rezultat}
\begin{figure}[H]
	\centering
	%\includegraphics{cps_aliasing_1.jpg}
	\includegraphics[width=\linewidth]{tmp.jpg}
	\caption{Wykres rekonstrukcji w oparciu o funkcję sinc sygnału sinusoidalnego z wystąpieniem aliasingu dla częstotliwości $f_0$=100Hz oraz $f_s$=1 000Hz}
	\label{wykres dla eksperymentu 6}
\end{figure}
Obliczone parametry dla otrzymanego w rezuntacie sygnału ciągłego:
\begin{figure}[H]
	\centering
	%\includegraphics{cps_aliasing_1_miary.jpg}
	\includegraphics[width=\linewidth]{tmp.jpg}
	\caption{Oblicznone miary dla rekonstrukcji w oparciu o funkcję sinc sygnału sinusoidalnego z wystąpieniem aliasingu dla częstotliwości $f_0$=100Hz oraz $f_s$=1 000Hz} 
	\label{wartości dla eksperymentu 6}
\end{figure}
\begin{figure}[H]
	\centering
	%\includegraphics{cps_aliasing_2.jpg}
	\includegraphics[width=\linewidth]{tmp.jpg}
	\caption{Wykres rekonstrukcji w oparciu o funkcję sinc sygnału sinusoidalnego z wystąpieniem aliasingu dla częstotliwości $f_0$=220Hz oraz $f_s$=44 100Hz}
	\label{wykres dla eksperymentu 6.2}
\end{figure}
Obliczone parametry dla otrzymanego w rezuntacie sygnału ciągłego:
\begin{figure}[H]
	\centering
	%\includegraphics{cps_aliasing_2_miary.jpg}
	\includegraphics[width=\linewidth]{tmp.jpg}
	\caption{Oblicznone miary dla rekonstrukcji w oparciu o funkcję sinc sygnału sinusoidalnego z wystąpieniem aliasingu dla częstotliwości $f_0$=220Hz oraz $f_s$=44 100Hz} 
	\label{wartości dla eksperymentu 6.2}
\end{figure}
\begin{figure}[H]
	\centering
	%\includegraphics{cps_aliasing_3.jpg}
	\includegraphics[width=\linewidth]{tmp.jpg}
	\caption{Wykres rekonstrukcji w oparciu o funkcję sinc sygnału sinusoidalnego z wystąpieniem aliasingu dla częstotliwości $f_0$=440Hz oraz $f_s$=22 050Hz}
	\label{wykres dla eksperymentu 6.3}
\end{figure}
Obliczone parametry dla otrzymanego w rezuntacie sygnału ciągłego:
\begin{figure}[H]
	\centering
	%\includegraphics{cps_aliasing_3_miary.jpg}
	\includegraphics[width=\linewidth]{tmp.jpg}
	\caption{Oblicznone miary dla rekonstrukcji w oparciu o funkcję sinc sygnału sinusoidalnego z wystąpieniem aliasingu dla częstotliwości $f_0$=440Hz oraz $f_s$=22 050Hz}
	\label{wartości dla eksperymentu 6.3}
\end{figure}
%%%%%%%%%%%%%%%%%%%%%%%%%%%%%%%%%%%%%%%%%%%%%%%%%%%%%%%%%%%%%%%%%%%%%%%%%%%%%%%%%%%%%%%%%%%%%%%%%%%%%%%%%%%%%%%%%

%%%%%%%%%%%%%%%%%%%%%%%%%%%%%%%%%%%%%%%%%%%%%%%%%%%%%%%%%%%%%%%%%%%%%%%%%%%%%%%%%%%%%%%%%%%%%%%%%%%%%%%%%%%%%%%%%
\section{Wnioski}
\begin{itemize}
    \item Każdy sygnał ciągły można zdyskretyzować poprzez próbkowanie a następnie te próbki skwantyzować by ograniczyć dokładność zapisu oraz ilość pamięci potrzebnej do zapisu sygnału zdyskrekretyzowanego.
    \item Program poprawnie kwantyzuje oraz rekonstruuje sygnały.
    \item Program pozwala obliczenie 4 miar dla każdego skwantyzowanego sygnału.
    \item Przy odpowiednich warunkach może wystąpić zjawisko aliasingu.
\end{itemize}
 

%%%%%%%%%%%%%%%%%%%%%%%%%%%%%%%%%%%%%%%%%%%%%%%%%%%%%%%%%%%%%%%%%%%%%%%%%%%%%%%%%%%%%%%%%%%%%%%%%%%%%%%%%%%%%%%%%
% BIBLIOGRAFIA
%%%%%%%%%%%%%%%%%%%%%%%%%%%%%%%%%%%%%%%%%%%%%%%%%%%%%%%%%%%%%%%%%%%%%%%%%%%%%%%%%%%%%%%%%%%%%%%%%%%%%%%%%%%%%%%%%

\renewcommand\refname{Bibliografia}
\bibliographystyle{plain}


\bibliography{instrukcja1} Instrukcja do zadania 1 dostępna na platformie WIKAMP, na stronie porzedmiotu Cyfrowe Przetwarzanie sygnału.
\bibliography{instrukcja2} Instrukcja do zadania 2 dostępna na platformie WIKAMP, na stronie porzedmiotu Cyfrowe Przetwarzanie sygnału.

\end{document}
